\documentclass[11pt,a4paper]{article}

\usepackage[provide=*,spanish]{babel}
\usepackage[utf8]{inputenc}
\usepackage[T1]{fontenc}
\usepackage{lmodern}
\usepackage[margin=2.35cm]{geometry}
\usepackage{amsmath}
\usepackage{booktabs}
\usepackage{longtable}
\usepackage{tabularx}
\usepackage{array}
\usepackage{enumitem}
\usepackage{listings}
\usepackage{xcolor}
\usepackage{microtype}
\usepackage{parskip}
\usepackage{titlesec}
\usepackage{hyperref}

\definecolor{azulWCD}{RGB}{17,67,112}
\definecolor{azulClaro}{RGB}{232,241,249}
\definecolor{grisCodigo}{RGB}{247,247,247}
\definecolor{verdeCodigo}{RGB}{25,110,55}
\definecolor{rojoCodigo}{RGB}{145,35,35}

\hypersetup{
  colorlinks=true,
  linkcolor=azulWCD,
  urlcolor=azulWCD,
  citecolor=azulWCD,
  pdftitle={Informe técnico de G4WCDSteppingAction},
  pdfauthor={Jaime A. Betancourt}
}

\lstdefinestyle{cppstyle}{
  language=C++,
  backgroundcolor=\color{grisCodigo},
  basicstyle=\ttfamily\footnotesize,
  keywordstyle=\color{azulWCD}\bfseries,
  commentstyle=\color{verdeCodigo}\itshape,
  stringstyle=\color{rojoCodigo},
  numbers=left,
  numberstyle=\tiny\color{gray},
  numbersep=8pt,
  breaklines=true,
  breakatwhitespace=false,
  frame=single,
  rulecolor=\color{gray!40},
  tabsize=2,
  showstringspaces=false,
  xleftmargin=1em,
  framexleftmargin=1em,
  columns=fullflexible,
  keepspaces=true
}
\lstset{style=cppstyle}

\titleformat{\section}{\normalfont\Large\bfseries\color{azulWCD}}
  {\thesection.}{0.6em}{}
\titleformat{\subsection}{\normalfont\large\bfseries}
  {\thesubsection.}{0.5em}{}
\setlist{itemsep=0.25em,topsep=0.4em}
\renewcommand{\arraystretch}{1.18}
\setlength{\emergencystretch}{2em}

\newcommand{\codigo}[1]{\nolinkurl{#1}}
\newcommand{\archivo}[1]{\path{#1}}

\title{\textbf{Informe técnico}\\[0.4em]
\Large Mejoras de la clase \codigo{G4WCDSteppingAction}\\[0.35em]
\large Registro de la respuesta neutrónica y electromagnética del detector
Cherenkov de agua}
\author{Jaime A. Betancourt\\
Doctorado en Física, Universidad Industrial de Santander}
\date{10 de septiembre de 2026}

\begin{document}

\maketitle

\begin{abstract}
Este informe documenta la adaptación y reorganización de la clase \codigo{G4WCDSteppingAction}, integrada en la aplicación \codigo{G4WCDSimulator} del framework MEIGA. La clase registra, a nivel de paso, el transporte de neutrones y la respuesta secundaria producida en el detector Cherenkov de agua (WCD), incluyendo capturas de neutrones, interacciones inelásticas, fotones gamma, electrones, positrones y fotones ópticos Cherenkov. La versión mejorada normaliza las unidades de salida, incorpora identificadores
de corrida y evento, corrige la identificación del núcleo capturador, diferencia los datos de partículas primarias y secundarias, protege la escritura en ejecuciones multihilo y reemplaza archivos de texto ambiguos por tablas TSV con encabezados explícitos. El documento distingue la infraestructura preexistente de MEIGA de las contribuciones desarrolladas para el estudio del WCD con agua y soluciones de NaCl.
\end{abstract}

\tableofcontents
\newpage

\section{Objeto y alcance del informe}

\codigo{G4WCDSteppingAction} deriva de \codigo{G4UserSteppingAction}. Geant4 invoca automáticamente su método \codigo{UserSteppingAction(const G4Step* step)} después de cada paso de cada trayectoria transportada. Por esta razón, la clase constituye un punto apropiado para extraer información microscópica sobre posiciones, energías, procesos y
partículas secundarias.

El propósito de la adaptación no es modificar los modelos físicos internos de Geant4. Su función es \textbf{seleccionar, organizar y conservar la información
producida por dichos modelos} para estudiar la cadena de respuesta

\begin{equation}
 n \longrightarrow
 \text{interacción o captura} \longrightarrow
 \gamma \longrightarrow e^\pm \longrightarrow
 \gamma_{\mathrm{Ch}}.
\end{equation}

Esta trazabilidad permite relacionar el transporte de neutrones con la producción electromagnética y óptica que finalmente puede contribuir a la señal del WCD. El código analizado corresponde a los archivos \archivo{G4WCDSteppingAction-mejorado.cc} y \archivo{G4WCDSteppingAction-mejorado.h}.

\section{Diferenciación entre MEIGA y la contribución de la tesis}

MEIGA es un framework desarrollado previamente para construir y ejecutar
simulaciones de detectores mediante Geant4. Por tanto, no constituyen aportes
originales de esta investigación:

\begin{itemize}
  \item la arquitectura general del framework;
  \item las clases base y el ciclo de ejecución de Geant4;
  \item la interfaz \codigo{G4UserSteppingAction};
  \item los mecanismos generales de transporte de partículas;
  \item los modelos físicos suministrados por Geant4.
\end{itemize}

Sobre esa infraestructura, el trabajo de tesis aporta la especialización de la
acción de paso para el estudio del WCD y, en particular:

\begin{enumerate}
  \item la selección diferenciada de neutrones incidentes, pasos de neutrones e   interacciones hadrónicas relevantes;
  \item el registro específico de capturas de neutrones y de sus productos secundarios;
  \item la consulta del núcleo objetivo de la interacción mediante sus números  atómico y másico, $Z$ y $A$;
  \item el seguimiento de la respuesta electromagnética secundaria mediante fotones gamma, electrones y positrones;
  \item la identificación de fotones ópticos creados por el proceso \codigo{Cerenkov};
  \item la definición de archivos tabulados con columnas, unidades e identificadores inequívocos;
  \item la incorporación de escritura protegida para ejecuciones multihilo;
  \item una estructura de datos adecuada para calcular eficiencias, espectros y distribuciones espaciales fuera de Geant4.
\end{enumerate}

En consecuencia, la contribución no corresponde a la creación de MEIGA ni de los procesos físicos de Geant4, sino a la \textbf{adaptación funcional del framework y al desarrollo del sistema de adquisición de datos simulados} requerido para caracterizar el WCD con agua pura y agua dopada con NaCl.

\section{Interfaz de la clase}

La cabecera conserva la firma del constructor empleada por \codigo{G4WCDSimulator}:

\begin{lstlisting}
G4WCDSteppingAction(const G4WCDConstruction* detector,
                    G4WCDEventAction* eventAction,
                    Event& event);

~G4WCDSteppingAction() override = default;

void UserSteppingAction(const G4Step* step) override;
\end{lstlisting}

Las referencias a la construcción del detector, a la acción de evento y al
objeto \codigo{Event} se mantienen para preservar la interfaz de MEIGA y permitir
extensiones posteriores. Se eliminan miembros que no contribuían al registro
activo, entre ellos el mapa \codigo{fReactionsPerEvent} y variables auxiliares
sin uso en esta implementación. El destructor se declara
\codigo{= default}, lo que expresa que la clase no administra manualmente
recursos cuya liberación requiera lógica adicional.

\section{Arquitectura de la implementación}

La implementación se divide en dos niveles:

\begin{enumerate}
  \item un espacio de nombres anónimo que contiene constantes y funciones
  auxiliares de uso interno;
  \item el método \codigo{UserSteppingAction}, que clasifica cada paso según la
  partícula y el proceso físico.
\end{enumerate}

Esta separación evita repetir bloques de apertura, escritura y formateo de
archivos. Las funciones auxiliares principales son:

\begin{description}
  \item[\codigo{ProcessName}] devuelve el nombre de un proceso y responde
  \codigo{none} cuando el puntero es nulo.
  \item[\codigo{MaterialName}] devuelve el nombre del material del punto de
  paso y maneja de forma segura la ausencia de material.
  \item[\codigo{CurrentRunId} y \codigo{CurrentEventId}] recuperan los
  identificadores de la corrida y del evento activos.
  \item[\codigo{AppendTsv}] añade una fila a un archivo TSV, crea el encabezado
  cuando el archivo está vacío y protege la operación con un mutex.
  \item[\codigo{StepRow}] construye una fila normalizada para una trayectoria.
  \item[\codigo{WriteSecondaries}] registra las partículas creadas en el paso.
  \item[\codigo{WriteCapture}] registra la captura y el núcleo objetivo.
  \item[\codigo{IsGammaInteraction} e \codigo{IsNeutronInteraction}]
  centralizan los procesos aceptados por los filtros.
  \item[\codigo{WriteCerenkovPhotons}] selecciona los fotones ópticos cuyo
  proceso creador es \codigo{Cerenkov}.
\end{description}

\section{Modelo común de los datos de salida}

\subsection{Unidades explícitas}

Geant4 trabaja internamente con un sistema coherente de unidades. La versión
anterior mezclaba valores sin conversión con valores divididos por
\codigo{MeV}, lo que hacía ambiguas algunas columnas. La nueva versión efectúa
la conversión en el punto de escritura:

\begin{lstlisting}
prePosition.x() / cm
pre->GetKineticEnergy() / MeV
track->GetGlobalTime() / ns
\end{lstlisting}

Por tanto, todas las posiciones se expresan en \textbf{cm}, las energías
cinéticas en \textbf{MeV} y los tiempos en \textbf{ns}.

\subsection{Identificación de corrida, evento y trayectoria}

Cada fila incorpora \codigo{run_id} y \codigo{event_id}. Además se registran
\codigo{track_id}, \codigo{parent_id} y \codigo{step_number}. Esta combinación
permite reconstruir la pertenencia de cada interacción a una historia concreta
y evita utilizar el número de fila como identificador físico.

Para una partícula primaria, \codigo{parent_id} es cero y
\codigo{creator_process} se registra como \codigo{none}. Para una secundaria,
el proceso creador permite establecer su procedencia física.

\subsection{Encabezados y formato TSV}

Los archivos utilizan valores separados por tabuladores y extensión
\codigo{.tsv}. \codigo{AppendTsv} escribe el encabezado solamente si el archivo
no existe o está vacío:

\begin{lstlisting}
if (needsHeader) {
  output << header << '\n';
}
output << row << '\n';
\end{lstlisting}

El encabezado comienza con el carácter \codigo{\#}; de esta forma puede ser
tratado como comentario por herramientas de análisis. Se definen tres esquemas:
pasos, capturas y partículas secundarias.

\subsection{Esquema de pasos}

\small
\begin{longtable}{>{\ttfamily}p{0.27\textwidth}p{0.64\textwidth}}
\toprule
\normalfont\textbf{Columna} & \textbf{Significado}\\
\midrule
\endfirsthead
\toprule
\normalfont\textbf{Columna} & \textbf{Significado}\\
\midrule
\endhead
run\_id & Identificador de la corrida actual.\\
event\_id & Identificador del evento actual.\\
particle & Nombre de la partícula transportada.\\
track\_id & Identificador de su trayectoria.\\
parent\_id & Identificador de la trayectoria progenitora.\\
step\_number & Número de paso dentro de la trayectoria.\\
creator\_process & Proceso que creó la trayectoria; \codigo{none} para primarias.\\
step\_process & Proceso que definió el final del paso.\\
material & Material en el punto anterior al paso.\\
pre\_x\_cm,\newline pre\_y\_cm,\newline pre\_z\_cm & Posición inicial del paso, en cm.\\
post\_x\_cm,\newline post\_y\_cm,\newline post\_z\_cm & Posición final del paso, en cm.\\
pre\_energy\_MeV & Energía cinética anterior al paso, en MeV.\\
post\_energy\_MeV & Energía cinética posterior al paso, en MeV.\\
global\_time\_ns & Tiempo global de la trayectoria, en ns.\\
\bottomrule
\end{longtable}
\normalsize

\subsection{Esquemas de capturas y secundarias}

El esquema de capturas añade \codigo{target_Z} y \codigo{target_A}. El esquema
de secundarias incluye la categoría de la interacción, la identidad de la
partícula padre, el número de paso del padre y
\codigo{secondary_index}. Este último corresponde a la posición de la
secundaria dentro del vector generado en el paso.

El identificador definitivo de trayectoria de una secundaria puede no estar
asignado todavía cuando se consulta \codigo{GetSecondaryInCurrentStep()}. Por
ello, \codigo{track_id} puede ser cero en esta etapa. Para identificar el
registro deben conservarse conjuntamente corrida, evento, trayectoria del
padre, paso del padre e índice de secundaria.

\section{Escritura segura de resultados}

\codigo{AppendTsv} usa un mutex global al módulo y un bloqueo RAII:

\begin{lstlisting}
G4Mutex outputMutex = G4MUTEX_INITIALIZER;

void AppendTsv(const G4String& fileName,
               const char* header,
               const G4String& row)
{
  G4AutoLock lock(&outputMutex);
  // Apertura, encabezado y escritura de la fila.
}
\end{lstlisting}

El bloqueo impide que dos hilos mezclen parcialmente sus filas en un mismo
archivo. También protege la verificación y creación del encabezado. Si el
archivo no puede abrirse, la clase genera una excepción fatal con un mensaje
que solicita comprobar la existencia y los permisos de
\archivo{Datos-simulacion/}.

Esta estrategia prioriza la integridad y simplicidad de los datos. La apertura
y cierre por fila, junto con la serialización de las escrituras, puede tener un
costo apreciable en corridas extensas. Si el rendimiento de entrada/salida se
convierte en una limitación, podrá evaluarse el uso de archivos por hilo o del
sistema de análisis de Geant4.

\section{Lógica de \texorpdfstring{\codigo{UserSteppingAction}}{UserSteppingAction}}

\subsection{Validación inicial}

Antes de acceder a los objetos del paso se comprueba que existan:

\begin{lstlisting}
if (!step || !step->GetTrack() || !step->GetPreStepPoint() ||
    !step->GetPostStepPoint()) {
  return;
}
\end{lstlisting}

El nombre del proceso se obtiene mediante \codigo{ProcessName}, evitando
desreferenciar directamente un posible puntero nulo retornado por
\codigo{GetProcessDefinedStep()}.

\subsection{Historia completa}

Cada paso se escribe inicialmente en \archivo{pasos-particulas.tsv}. Este
archivo proporciona la historia más detallada, pero también puede crecer con
rapidez. La propia implementación indica que su llamada puede comentarse cuando
solo se requieran productos filtrados.

\subsection{Selección de neutrones incidentes}

Un neutrón se registra como incidente únicamente cuando es primario y se
encuentra en su primer paso:

\begin{lstlisting}
if (track->GetParentID() == 0 &&
    track->GetCurrentStepNumber() == 1) {
  AppendTsv("neutrones-incidentes.tsv", stepHeader, row);
}
\end{lstlisting}

Esta condición corrige el criterio anterior, que seleccionaba cualquier
neutrón en su primer paso y podía incluir neutrones secundarios como si fueran
partículas inyectadas.

\subsection{Interacciones de neutrones}

\codigo{IsNeutronInteraction} conserva los procesos relevantes para este
análisis:

\begin{itemize}
  \item \codigo{hadElastic}: dispersión elástica y moderación;
  \item \codigo{nCapture}: captura neutrónica;
  \item \codigo{neutronInelastic}: interacción inelástica.
\end{itemize}

Todos sus pasos se escriben en
\archivo{interacciones-neutrones.tsv}. La historia completa de cada neutrón,
incluidos los pasos de transporte, se conserva por separado en
\archivo{pasos-neutrones.tsv}.

\section{Captura de neutrones y núcleo objetivo}

La versión anterior intentaba inferir el capturador tomando \codigo{material->GetElement(0)}. Ese método devuelve el primer elemento de la definición del material, no necesariamente el núcleo seleccionado por la interacción.
%Además, una condición correspondiente a oxígeno ($Z=8$) escribía en un archivo cuyo nombre hacía referencia a cloro-35.

La versión mejorada consulta el proceso hadrónico que definió el paso:

\begin{lstlisting}
const auto* process = post->GetProcessDefinedStep();
const auto* hadronicProcess =
  dynamic_cast<const G4HadronicProcess*>(process);

if (hadronicProcess && hadronicProcess->GetTargetNucleus()) {
  targetZ = hadronicProcess->GetTargetNucleus()->GetZ_asInt();
  targetA = hadronicProcess->GetTargetNucleus()->GetA_asInt();
}
\end{lstlisting}

Así, una captura en $^{1}\mathrm{H}$ puede representarse mediante
$(Z,A)=(1,1)$; una captura en $^{35}\mathrm{Cl}$, mediante $(17,35)$; y una
captura en $^{37}\mathrm{Cl}$, mediante $(17,37)$. Si el proceso no proporciona
un núcleo objetivo, ambos campos conservan el valor centinela $-1$.

Para cada \codigo{nCapture} se escribe una fila en
\archivo{capturas-neutrones.tsv}. Sus productos se recorren mediante
\codigo{GetSecondaryInCurrentStep()} y se almacenan en
\archivo{secundarias-captura-neutron.tsv}. Esta separación permite construir
distribuciones por isótopo capturador y espectros de los fotones gamma de
desexcitación.

\section{Interacciones gamma y partículas cargadas}

Todos los pasos de fotones gamma se escriben en
\archivo{pasos-gamma.tsv}. Se consideran de interés específico los procesos:

\begin{center}
\begin{tabular}{ll}
\toprule
\textbf{Nombre de Geant4} & \textbf{Proceso físico}\\
\midrule
\codigo{phot} & efecto fotoeléctrico\\
\codigo{compt} & dispersión Compton\\
\codigo{conv} & producción de pares\\
\codigo{Rayl} & dispersión Rayleigh\\
\bottomrule
\end{tabular}
\end{center}

La información se distribuye en tres productos complementarios:
\begin{itemize}
  \item \archivo{interacciones-gamma.tsv}, para los pasos seleccionados;
  \item \archivo{secundarias-interaccion-gamma.tsv}, para sus productos;
  \item \archivo{pasos-electrones.tsv}, para los pasos de electrones y
  positrones.
\end{itemize}
En conjunto, estos archivos permiten seguir la transferencia de energía desde
los gamma de captura hacia las partículas cargadas que pueden producir luz
Cherenkov.

\section{Identificación de fotones Cherenkov}

La versión anterior verificaba si el proceso que limitó el paso del electrón
era \codigo{Cerenkov}. Esta condición puede omitir producción óptica, porque
Cherenkov genera fotones secundarios a lo largo de un paso sin ser
necesariamente el proceso que determina su longitud.

La versión mejorada recorre las secundarias de cada paso y exige simultáneamente
que la partícula sea un fotón óptico y que su proceso creador sea
\codigo{Cerenkov}:

\begin{lstlisting}
if (secondary->GetParticleDefinition() !=
      G4OpticalPhoton::Definition() ||
    ProcessName(creator) != "Cerenkov") {
  continue;
}
\end{lstlisting}

El filtro se aplica a cualquier partícula padre, no solo a electrones y
positrones. Por tanto, también permite registrar luz producida por otras
partículas cargadas presentes en el fondo cósmico. Los resultados se guardan
en \archivo{fotones-cherenkov.tsv}.

Este archivo cuantifica \textbf{fotones Cherenkov producidos}; no representa
por sí mismo los fotones que alcanzan el fotocátodo ni los fotoelectrones
generados por el PMT. La eficiencia instrumental requiere aplicar además el
transporte óptico, la geometría del fotocátodo, la eficiencia cuántica y el
criterio de disparo.

\section{Archivos generados}

Todos los archivos se escriben dentro de
\archivo{Datos-simulacion/}. La carpeta debe existir y disponer de permisos de
escritura antes de iniciar el programa.

\small
\begin{longtable}{>{\raggedright\arraybackslash}p{0.32\textwidth}
                  >{\raggedright\arraybackslash}p{0.23\textwidth}
                  >{\raggedright\arraybackslash}p{0.36\textwidth}}
\toprule
\textbf{Archivo} & \textbf{Condición} & \textbf{Uso principal}\\
\midrule
\endfirsthead
\toprule
\textbf{Archivo} & \textbf{Condición} & \textbf{Uso principal}\\
\midrule
\endhead
\archivo{pasos-particulas.tsv} & Todo paso & Diagnóstico completo de la simulación.\\
\archivo{pasos-neutrones.tsv} & Partícula neutrón & Transporte y moderación de neutrones.\\
\archivo{neutrones-incidentes.tsv} & Neutrón primario, paso 1 & Número y condiciones iniciales de los neutrones inyectados.\\
\archivo{interacciones-neutrones.tsv} & \codigo{hadElastic}, \codigo{nCapture} o \codigo{neutronInelastic} & Conteo, energía y posición de interacciones neutrónicas.\\
\archivo{capturas-neutrones.tsv} & \codigo{nCapture} & Capturas con material y núcleo objetivo.\\
\archivo{secundarias-captura-neutron.tsv} & Secundarias de \codigo{nCapture} & Productos de captura y espectros gamma.\\
\texttt{secundarias-}\newline\texttt{inelastica-neutron.tsv} & Secundarias de \codigo{neutronInelastic} & Productos de interacciones inelásticas.\\
\archivo{pasos-gamma.tsv} & Partícula gamma & Transporte de fotones gamma.\\
\archivo{interacciones-gamma.tsv} & \codigo{phot}, \codigo{compt}, \codigo{conv} o \codigo{Rayl} & Procesos electromagnéticos seleccionados.\\
\texttt{secundarias-}\newline\texttt{interaccion-gamma.tsv} & Secundarias de interacción gamma & Electrones, positrones y otros productos.\\
\archivo{pasos-electrones.tsv} & Partícula $e^-$ o $e^+$ & Transporte de partículas cargadas.\\
\archivo{fotones-cherenkov.tsv} & Fotón óptico creado por \codigo{Cerenkov} & Producción de luz Cherenkov.\\
\bottomrule
\end{longtable}
\normalsize

Los archivos se abren en modo \emph{append}. Por ello, una ejecución posterior
añade filas a los resultados existentes. Aunque \codigo{run_id} y
\codigo{event_id} separan registros durante una ejecución, esos identificadores
pueden reiniciarse al lanzar nuevamente el programa. Para evitar mezclar
campañas independientes se recomienda vaciar, renombrar o archivar la carpeta
de salida antes de cada conjunto de simulaciones.

\section{Interpretación de resultados}

\subsection{Eficiencia de captura}

Si cada fila de \archivo{neutrones-incidentes.tsv} corresponde a un neutrón
primario y cada fila de \archivo{capturas-neutrones.tsv} a una captura, puede
calcularse

\begin{equation}
\eta_{\mathrm{cap}}=
\frac{N_{\mathrm{cap}}}{N_{\mathrm{inc}}}.
\end{equation}

El conteo puede desagregarse por energía incidente, material y núcleo objetivo.
Para comparar agua pura con soluciones de NaCl debe mantenerse el mismo número
de neutrones incidentes, la misma geometría, la misma fuente y la misma lista de
física.

\subsection{Espectros y distribuciones espaciales}

Los espectros de secundarios se construyen a partir de
\codigo{kinetic_energy_MeV}. Las distribuciones espaciales de interacciones y
capturas se obtienen con las coordenadas en cm. Los análisis por evento deben
agrupar primero por \codigo{run_id} y \codigo{event_id} para impedir que varios
pasos de la misma cascada se interpreten como eventos independientes.

\subsection{De la captura a la eficiencia del detector}

La eficiencia de captura no equivale a la eficiencia de detección. Una
definición instrumental debe incluir un criterio observable, por ejemplo,

\begin{equation}
\eta_{\mathrm{det}}(E_n)=
\frac{N_{\text{eventos que superan el umbral de señal}}(E_n)}
     {N_{\text{neutrones incidentes}}(E_n)}.
\end{equation}

El numerador puede depender del número de fotoelectrones, de la ventana
temporal, del umbral electrónico y de la respuesta del PMT. Las tablas
producidas por esta clase describen etapas físicas de la cadena, pero deben
combinarse con el registro del fotocátodo y de la electrónica simulada para
obtener la eficiencia final del WCD.

\section{Síntesis comparativa de las mejoras}

\small
\begin{longtable}{>{\raggedright\arraybackslash}p{0.27\textwidth}
                  >{\raggedright\arraybackslash}p{0.30\textwidth}
                  >{\raggedright\arraybackslash}p{0.34\textwidth}}
\toprule
\textbf{Aspecto} & \textbf{Situación anterior} & \textbf{Versión mejorada}\\
\midrule
\endfirsthead
\toprule
\textbf{Aspecto} & \textbf{Situación anterior} & \textbf{Versión mejorada}\\
\midrule
\endhead
Unidades & Mezcla de unidades internas y valores convertidos. & Conversión explícita a cm, MeV y ns.\\
Encabezados & Archivos sin nombres de columnas. & Encabezados comentados y esquemas coherentes.\\
Identidad del evento & No se escribían corrida ni evento. & Inclusión de \codigo{run_id} y \codigo{event_id}.\\
Neutrón incidente & Todo neutrón en su primer paso. & Solo neutrón primario con \codigo{parent_id=0} y paso 1.\\
Núcleo capturador & Primer elemento del material; archivo de cloro asociado a una condición de oxígeno. & Núcleo objetivo del proceso hadrónico mediante $Z$ y $A$.\\
Secundarias & Identificadores del padre reutilizados de forma ambigua. & Datos separados del padre y la secundaria.\\
Cherenkov & Comparación con el proceso que limitó el paso del electrón. & Verificación del tipo de partícula y del proceso creador del fotón óptico.\\
Punteros a procesos & Acceso directo sin comprobación. & Manejo seguro con respuesta \codigo{none}.\\
Multihilo & Escrituras potencialmente concurrentes. & Protección de cada fila con \codigo{G4AutoLock}.\\
Duplicación & Bloques repetidos y archivos parcialmente redundantes. & Funciones auxiliares y archivos con propósitos diferenciados.\\
Mapa de reacciones & Estructura inactiva y limpiada por paso. & Eliminación del estado que no acumulaba información útil.\\
\bottomrule
\end{longtable}
\normalsize

\section{Condiciones de uso y limitaciones}

\begin{enumerate}
  \item \textbf{Directorio de salida.} Debe existir
  \archivo{Datos-simulacion/}; de lo contrario, la ejecución termina mediante
  la excepción \codigo{WCD-OUTPUT-001}.
  \item \textbf{Volumen de datos.} Registrar todos los pasos puede producir
  archivos grandes. Para corridas de producción puede deshabilitarse
  \archivo{pasos-particulas.tsv} si no se requiere diagnóstico completo.
  \item \textbf{Orden multihilo.} El bloqueo garantiza filas íntegras, pero el
  orden global de las filas no tiene por qué coincidir con el orden numérico de
  los eventos. El análisis debe usar los identificadores escritos.
  \item \textbf{Identidad de secundarias.} El \codigo{track_id} puede no estar
  asignado todavía al consultar las secundarias del paso.
  \item \textbf{Alcance óptico.} El conteo de
  \archivo{fotones-cherenkov.tsv} no sustituye el conteo de fotones detectados o
  fotoelectrones.
  \item \textbf{Validación.} El código y este informe fueron contrastados
  estáticamente. La compilación y la prueba de ejecución deben realizarse con
  el repositorio completo, Geant4 10.07.p04, la lista de física y los conjuntos
  de datos nucleares empleados en la investigación.
\end{enumerate}

\section{Integración y verificación recomendadas}

Los archivos deben instalarse con los nombres esperados por el proyecto:

\begin{center}
\begin{tabular}{ll}
\toprule
\textbf{Archivo de trabajo} & \textbf{Nombre dentro del repositorio}\\
\midrule
\archivo{G4WCDSteppingAction-mejorado.cc} & \archivo{G4WCDSteppingAction.cc}\\
\archivo{G4WCDSteppingAction-mejorado.h} & \archivo{G4WCDSteppingAction.h}\\
\bottomrule
\end{tabular}
\end{center}

La ubicación de destino es
\archivo{src/Applications/G4WCDSimulator/}. Después de reemplazarlos se
recomienda:

\begin{enumerate}
  \item comprobar los cambios con \codigo{git diff --check};
  \item configurar y compilar el proyecto con Geant4 10.07.p04;
  \item ejecutar una corrida corta;
  \item verificar encabezados, unidades, materiales y nombres de procesos;
  \item comprobar manualmente casos conocidos de captura en hidrógeno y cloro;
  \item comparar los conteos con los obtenidos por la versión anterior;
  \item ejecutar finalmente las campañas de producción.
\end{enumerate}

\section{Conclusiones}

La versión mejorada de \codigo{G4WCDSteppingAction} transforma un registro
heterogéneo y parcialmente redundante en un sistema tabulado, trazable y más
seguro. Las correcciones más relevantes son la normalización de unidades, la
identificación del núcleo capturador a partir del proceso hadrónico, la
separación correcta entre primarias y secundarias y el reconocimiento de la luz
Cherenkov mediante el proceso creador del fotón óptico.

Desde la perspectiva de la tesis, el aporte consiste en adaptar MEIGA para
extraer observables específicamente relacionados con la respuesta de un WCD a
neutrones en agua pura y soluciones de NaCl. Esta capacidad conecta la física
de interacción neutrónica con la generación de partículas secundarias y luz
Cherenkov, y proporciona una base reproducible para los análisis de captura,
respuesta espectral y eficiencia del detector.

La integración definitiva queda condicionada a la compilación y validación con
Geant4 10.07.p04. Esta condición debe conservarse explícitamente en el
repositorio hasta completar las pruebas correspondientes.

\end{document}
